\newpage
\section{Sieci kwantowe}

Obliczenia kwantowe oraz kwantowe przetwarzanie danych są relatywnie nowymi dyscyplinami, w~których zasady fizyki kwantowej są wykorzystywane do przechowywanie i~przetwarzania danych. \cite{Nakahara2013}
Uczenie maszynowe przyczyniło się do rewolucyjnych zastosowań zarówno w~społeczeństwie jak i~w świecie nauki, jednakże wraz z~ciągle słabnącym prawem Moore’a poszukiwane są nowe ścieżki celem podtrzymania dynamizmu rozwoju technologicznego. Jedną z~tych ścieżek jest kwantowe uczenie maszynowe, które podąża za rozwojem pierwszych urządzeń kwantowych, których powstanie jest skutkiem próby ominięcia ograniczeń technologicznych i~termodynamicznych klasycznej komputacji.\cite{training_deep_quantum_nn}
O ile obecnie wykorzystywane technologie podlegają prawom fizyki kwantowej o~tyle logika pracy opiera się na fizyce klasycznej. 
Obliczenia kwantowe zasadnicza różnią się dwoma kluczowymi aspektami:
\begin{itemize}
    \item Superpozycja stanu. Układ może znajdować się w~dwóch stanach jednocześnie, dopóki nie zostanie wykonany jego pomiar.
    \item Splątanie stanów. Pomiar jednego układu wpływa na stan innych z~nim splątanych układów.
\end{itemize}

Wykorzystanie kwantowego podejścia do opisu danych umożliwia znaczące przyspieszenie obliczeń oraz m.in. systemy kryptograficzne odporne na niektóre typy ataków. \cite{Nakahara2013}

\subsection{Czym jest qubit?}\cite{quantum_intro}

Qubit (ang. \textit{quantum bit}) jest podstawową jednostką informacji w~obliczeniach kwantowych, odpowiednikiem bitu w~obliczeniach klasycznych. W~przeciwieństwie do bitu, który może przyjmować jedną z~dwóch wartości (0 lub 1), qubit może znajdować się w~stanie superpozycji obu tych wartości jednocześnie. Stany qubitu można zwizualizować na sferze Blocha (o promieniu równym 1) w~układzie kartezjańskim. Zgodnie z~konwencją: oś z~jest skierowana w~górę a~oś y~w~bok.
Szczególnymi przypadkami stanów qubitu są:

\begin{table}[H]
\caption{Reprezentacja podstawowych stanów qubitów na sferze Blocha.}
\label{tab:qubit_bloch}
\centering
\resizebox{\textwidth}{!}{
\begin{tabular}{l c c c}
\hline
\textbf{Oznaczenie} & \textbf{Stan} & \textbf{Współrzędne na sferze Blocha} \\
\hline
$\ket{0}$ & $\ket{0}$ & (0, 0, 1) \\
$\ket{1}$ & $\ket{1}$ & (0, 0, -1) \\
$\ket{+}$ & $\frac{1}{\sqrt{2}} \left( \ket{0}+\ket{1}\right)$ & (1, 0, 0) \\
$\ket{-}$ & $\frac{1}{\sqrt{2}} \left( \ket{0}-\ket{1}\right)$ & (-1, 0, 0) \\
$\ket{i}$ & $\frac{1}{\sqrt{2}} \left( \ket{0}+i\ket{1}\right)$ & (0, 1, 0) \\
$\ket{-i}$ & $\frac{1}{\sqrt{2}} \left( \ket{0}-i\ket{1}\right)$ & (0, -1, 0) \\
\hline
\end{tabular}
}
\end{table}

\begin{figure}[H]
    \centering
    \begin{minipage}[b]{0.48\textwidth}
        \centering
        \includegraphics[width=\linewidth,height=6cm,keepaspectratio]{Figures/bloch_from_introduction_pg_76.png}
        \caption{Sfera Blocha z~zaznaczonymi podstawowymi stanami qubitów. \cite{quantum_intro}}
        \label{fig:bloch_sphere}
    \end{minipage}
    \hfill
    \begin{minipage}[b]{0.48\textwidth}
        \centering
        \includegraphics[width=\linewidth,height=6cm,keepaspectratio]{Figures/bloch_angles.png}
    \caption{Sfera Blocha z~kątami $\theta$ oraz $\phi$ \cite{quantum_intro}.}
    \label{fig:bloch_angles}
    \end{minipage}
\end{figure}


Interpetowalnym wynikiem pomiaru qubitu jest zawsze jedna z~dwóch wartości: 0 lub 1. Prawdopodobieństwo uzyskania każdej z~tych wartości zależy od stanu, w~jakim znajduje się qubit przed pomiarem. Na przykład, jeśli qubit jest w~stanie $\ket{+}$, to istnieje 50\% szans na uzyskanie wyniku 0 i~50\% szans na uzyskanie wyniku 1 podczas pomiaru.

Prawdopodobieństwo wyznacza się ze wzoru:

\begin{equation}
\label{eq:quantum_state}
\ket{\psi} = \alpha \ket{0} + \beta \ket{1}, \quad
P(0) = |\braket{0|\psi}|^2 = |\alpha|^2, \quad
P(1) = |\braket{1|\psi}|^2 = |\beta|^2,
\end{equation}


gdzie $\ket{\Psi}$, to stan kwantowy, $P(0)$ to prawdopodobieństwo uzyskania wyniku 0, a~$P(1)$ to prawdopodobieństwo uzyskania wyniku 1. Wartości $\alpha$ i~$\beta$ są amplitudami prawdopodobieństwa dla stanów $\ket{0}$ i~$\ket{1}$ odpowiednio, przy czym muszą spełniać warunek normalizacji: $|\alpha|^2 + |\beta|^2 = 1$.

Jednocześnie wzór \ref{eq:quantum_state} można wyrazić:
\begin{equation}
    \ket{\psi} = \cos{\frac{\theta}{2}} \ket{0} + e^{i\phi}\sin{\frac{\theta}{2}} \ket{1}, \quad
\end{equation}

Gdzie $\theta$ to odległość zenitalna, czyli miara kąta między wektorem $\vec{OM}$ ($O$ - początek układu współrzędnych sfery Blocha, $M$ punkt na sferze Blocha) a~dodatnią półosią $OZ$, natomiast $\phi$ to długość azymutalna, czyli miara kąta między rzutem prostokątnym wektora $\vec{OM}$ na płaszczyznę OXY a~dodatnią półosią OX. (vide rys. \ref{fig:bloch_angles}). Właściwości tych miar kątowych uzasadniają skalowanie zbioru danych do zakresu $[-1,1]$, które umożliwia wprowadzenie danych diagnostycznych do stanu kwantowego jako kąty wyrażone w~radianach.

Pomiary można wykonywać w~różnych bazach (względem różnych osi na sferze Blocha), co pozwala na uzyskanie różnych informacji o~stanie qubitu przed pomiarem.

\subsection{Bramki kwantowe}\cite{quantum_intro}

Bramki kwantowe stanowią odpowiednik klasycznych bramek logicznych w~obliczeniach kwantowych. Są transformacją, liniowym mapowaniem stanu qubitu w~inne stany. Po przekształceniu całkowite prawdopodobieństwo stanu wciąż wynosi 1. W~ujęciu algebry liniowej bramki kwantowe są macierzami. Poniższa tabela przedstawia wybrane bramki kwantowe wraz z~ich działaniem na bazie, reprezentacją macierzową oraz interpretacją funkcjonalną.
Każda bramka kwantowa jest macierzą unitarną, co oznacza:
\[
U^{\dagger} U~= U~U^{\dagger} = I,
\]
gdzie $U$ jest macierzą bramki kwantowej, $U^{\dagger}$ jest jej sprzężeniem hermitowskim (transpozycja macierzy + sprzężenie), a~$I$ jest macierzą jednostkową.


\begin{table}[H]
\caption{Przykładowe bramki kwantowe działające na pojedynczym qubicie.}
\label{tab:bramki_kwantowe}
\resizebox{\textwidth}{!}
{
    \centering
    \begin{tabular}{l l l l}
        \hline
        \textbf{Bramka} & \textbf{Działanie} & \textbf{Forma macierzowa} & \textbf{Interpretacja} \\
        \hline
        Identity & $\begin{aligned} I\ket{0} &= \ket{0} \\ I\ket{1} &= \ket{1} \end{aligned}$ & $I = \begin{pmatrix} 1 & 0 \\ 0 & 1 \end{pmatrix}$ & Brak zmian (operacja neutralna). \\
        \hline
        Pauli $X$ & $\begin{aligned} X\ket{0} &= \ket{1} \\ X\ket{1} &= \ket{0} \end{aligned}$ & $X = \begin{pmatrix} 0 & 1 \\ 1 & 0 \end{pmatrix}$ & Obrót o~$\pi$ wokół osi $x$. \\
        \hline
        Pauli $Y$ & $\begin{aligned} Y\ket{0} &= i\ket{1} \\ Y\ket{1} &= -i\ket{0} \end{aligned}$ & $Y = \begin{pmatrix} 0 & -i \\ i~& 0 \end{pmatrix}$ & Obrót o~$\pi$ wokół osi $y$ \\
        \hline
        Pauli $Z$ & $\begin{aligned} Z\ket{0} &= \ket{0} \\ Z\ket{1} &= -\ket{1} \end{aligned}$ & $Z = \begin{pmatrix} 1 & 0 \\ 0 & -1 \end{pmatrix}$ & Obrót o~$\pi$ wokół osi $z$. \\
        \hline
        Phase $S$ & $\begin{aligned} S\ket{0} &= \ket{0} \\ S\ket{1} &= i\ket{1} \end{aligned}$ & $S = \begin{pmatrix} 1 & 0 \\ 0 & i~\end{pmatrix}$ & Obrót o~$\frac{\pi}{2}$ wokół osi $z$. \\
        \hline
        $T$ & $\begin{aligned} T\ket{0} &= \ket{0} \\ T\ket{1} &= e^{i\pi/4}\ket{1} \end{aligned}$ & $T = \begin{pmatrix} 1 & 0 \\ 0 & e^{i\pi/4} \end{pmatrix}$ & Obrót o~$\frac{\pi}{4}$ wokół osi $z$. \\
        \hline
        Hadamard $H$ & $\begin{aligned} H\ket{0} &= \frac{\ket{0} + \ket{1}}{\sqrt{2}} \\ H\ket{1} &= \frac{\ket{0} - \ket{1}}{\sqrt{2}} \end{aligned}$ & $H = \frac{1}{\sqrt{2}} \begin{pmatrix} 1 & 1 \\ 1 & -1 \end{pmatrix}$ & Obrót od $\pi$ jednocześnie wokół oś $x$ i~$z$ \\
        \hline
    \end{tabular}
}

\end{table}

Szczególną własnością powyższych bramek jest:
\[X^2 = Y^2 = Z^2 = S^4 = T^8 = I,\]

\subsection{Architektura hybrydowych kwantowo-klasycznych sieci neuronowych}

W architekturze głębokich sieci kwantowych jako analog perceptronu autorzy \cite{training_deep_quantum_nn} proponują dowolny unitarny operator z~$m$ qubitami na wejściu i~$n$ na wyjściu. Operator ten jest zależny od $(2^{m+n})^2-1$ parametrów. Qubity wejściowe są inicjalizowane nieznanym stanem $\rho^{in}$, a~qubity wyjściowe czystym stanem $\ket{0...0}_{out}$. Dla uproszczenia przyjmuje się $n = 1$. Model można zapisać w~następującej postaci:

\begin{equation}\label{eq:stan_wyjscie_1}
    \rho^{out} \equiv \operatorname{tr}_{in,hidden} \left[ U~(\rho^{in} \otimes \ket{0...0}_{hidden,out}\bra{0...0}) U^{\dagger} \right],
\end{equation}

gdzie:
\begin{itemize}
    \item $\operatorname{tr}_{in,hidden}$ oznacza częściowy ślad po qubitach wejściowych i~ukrytych.
    \item $U \equiv U^{out}U^L U^{L-1} ... U^1$
    \item $U^l$ jest iloczynem operatorów działających na warstwach l-1,l
\end{itemize}

Równoważnie można zapisać \ref{eq:stan_wyjscie_1} jako:

\begin{equation}\label{eq:stan_wyjscie_2}
    \rho^{out} \equiv \varepsilon^{out}(\varepsilon^L(...\varepsilon^2(\varepsilon^1(\rho^{in})))),
\end{equation}

gdzie:
\begin{itemize}
    \item $\varepsilon^l(X^{l-1}) = \operatorname{tr}_{l-1} \left(\prod_{j=m_l}^1 U_j^l (X^{l-1}\otimes\ket{0...0}_l\bra{0...0})\prod_{j=1}^{m_l}{U_j^l}^{\dagger}  \right)$
    \item $m_l$ - liczba perceptronów działających na warstwach $l-1$ oraz $l$,
    \item $j$ - indeks perceptronu w~warstwie $l$,
\end{itemize}

Uczona jest relacja między stanami kwantowymi. Dane treningowe są w~postaci par stanów kwantowych $(\ket{\phi_x^{in}},\ket{\phi_x^{out}})$, gdzie $x = 1,2,...,N$ to indeks przykładu treningowego z~$N$ próbek. Rozpatrywany jest scenariusz, w~którym $\ket{\phi_x^{out}}=V\ket{\phi_x^{in}}$, gdzie $V$ jest nieznanym unitarnym operatorem. 
Wprowadzona funkcja kosztu ocenia zbieżność $\rho_x^{out}$ z~$\ket{\phi_x^{out}}$:

\begin{equation}\label{eq:f_kosztu}
    C~= \frac{1}{N} \sum_{x=1}^{N} \bra{\phi_x^{out}} \rho_x^{out} \ket{\phi_x^{out}},
\end{equation}

Jest to uogólniona funkcja ryzyka wykorzystywana w~klasycznym uczeniu głębokich sieci kwantowych dla czystych stanów kwantowych. Przyjmuje ona wartości z~przedziału [0,1], gdzie 1 oznacza idealne dopasowanie. 
Sieć jest trenowana optymalizując funkcję kosztu \ref{eq:f_kosztu} poprzez dostosowanie parametrów bramek kwantowych w~sieci na każdym kroku iteracyjnym. 
\[U \rightarrow e^{i\epsilon K}U\]
gdzie:
\begin{itemize}
    \item $K$ - macierz parametrów operatora unitarnego. Dobrane początkowo tak, aby $C$ rosło jak najszybciej,
    \item $\epsilon$ - rozmiar kroku (ang. \textit{step size}).
\end{itemize}
do wyznaczenia $K$ wykorzystywane jest równanie:

\begin{equation}\label{eq:delta_C}
    \Delta C~= \frac{\epsilon}{N}\sum_{x=1}^{N}\sum_{l=1}^{L+1}\operatorname{tr}(\sigma_x^l\Delta\varepsilon^l(\rho_x^{l-1}))
\end{equation}

gdzie:
\begin{itemize}
    \item $L+1$ - warstwa wyjściowa,
    \item $\rho_x^{l-1}=\varepsilon^l(...\varepsilon^2(\varepsilon^1(\rho_x^{in})))$
    \item $\sigma_x^l =\mathcal{F}^{l+1}(...\mathcal{F}^L(\mathcal{F}^{out}(\ket{\phi^{out}}\bra{\phi}^{out}))...) $
    \item $\mathcal{F}(X) = \sum_{\alpha}A_{\alpha}^{\dagger}XA{\alpha}$ - sprzężny kanał kwantowy do $\varepsilon(X)=\sum_{\alpha}A_{\alpha}XA{\alpha}^{\dagger}$, czyli kanał przejścia z~$l+1$ do $l$,
\end{itemize}

Własnością tej struktury sieci jest fakt, że dla wyznaczenia $K_j^l$ dla warstwy $l-1$ oraz $l$, potrzebne są jedynie informacje z~tych dwóch warstw. Oznacza to, że aktualizacja lokalnych parametrów może być wykonywana lokalnie dla każdej warstwy, bez konieczności wyznaczania każdego z~parametrów sieci, jak ma to miejsce w~klasycznych sieciach neuronowych. $K$ zatem skaluje się z~szerokością sieci a~nie z~jej głębokością, co jest korzystne z~punktu widzenia efektywności obliczeniowej.

Ograniczeniem jest jednak ilość qubitów ze względu na ekspotencjalnie rosnącą przestrzeń Hilberta. Autorzy w~tym podejściu dostrzegają kluczową zaletę ograniczenie występowania tzw. zanikania gradientu (ang. \textit{barren plateau}) podczas trenowania sieci kwantowych, co jest powszechnym problemem w~klasycznych sieciach neuronowych oraz innych architekturach QNN.

\begin{tcolorbox}[title=Algorytm trenowania głębokiej sieci kwantowej,
    colback=blue!5!white,colframe=blue!75!black,
    breakable]
    
    \begin{enumerate}
        \item \textbf{Inicjalizacja} Krok $s=0$. Losowo inicjalizowane są parametry bramek kwantowych w~sieci.
        \item \textbf{Feedforward} Dla każdej pary stanów treningowych $\ket{\phi_x^{in}},\ket{\phi_x^{out}}$ wykonane są następujące działania:
        \begin{enumerate}
            \item Wyznaczony zostaje stan: \[\mathcal{T} =\rho_x^{l-1}(s)\otimes \ket{0...0}_l\bra{0...0} \]
            \item Zastosowane zostają operatory \[\mathcal{O} = \prod_{j=m_l}^{1}U_j^l(s)\mathcal{T}\prod_{j=1}^{m_l}{U_j^l(s)}^{\dagger} \]
            \item Wyznaczony zostaje ślad częściowy \[\rho_x^l(s) = \operatorname{tr}_{l-1}(\mathcal{O}) \]
            \item Przechowany zostaje stan $\rho_x^l(s)$ 
        \end{enumerate}
        Ten krok jest równoważny z~wyznaczeniem $\rho_x^{out}(s)$ zgodnie z~równaniem \ref{eq:stan_wyjscie_2}.
        \item \textbf{Aktualizacja parametrów}
        \begin{enumerate}
            \item Oblicza się funkcję kosztu $C(s)$ zgodnie z~równaniem \ref{eq:f_kosztu}.
            \item Wyznacza się każdą macierz $K_j^l(s)$ dla każdej warstwy $l$ i~każdego perceptronu $j$ w~tej warstwie:
            \begin{equation}
                K_j^l(s)=\eta \frac{2^{m_{l-1}}}{N}\sum_{x=1}^{N}\operatorname{tr}_{rest}M_j^l
            \end{equation}
            gdzie:
            \begin{itemize}
                \item $\eta$ - współczynnik uczenia,
                \item $\operatorname{tr}_{rest}$ - ślad częściowy po wszystkich qubitach oprócz tych, na których działa $U_j^l$,
                \item $M_j^l$ - komutator:
                \[\left[\prod_{\alpha=j}^{1}U_{\alpha}^l(\mathcal{T})\prod_{\alpha=1}^{j}{U_{\alpha}^l}^{\dagger},
                \prod_{\alpha=j+1}^{m_l}{U_{\alpha}^l}^{\dagger} (\mathbb{I}_{l-1}\otimes \sigma_x^l)\prod_{\alpha=m_l}^{j+1}U_{\alpha}^l\right]\]
            \end{itemize}
            \item \textbf{Aktualizacja operatorów unitarnych} Dla każdej warstwy l~i każdego perceptronu j~wykonane zostają operacje unitarne:
        \[U_j^l(s+1) = e^{i\epsilon K_j^l(s)}U_j^l(s)\]
            \item Zwiększ krok $s = s~+ \epsilon$ 
        \end{enumerate}
        \item \textbf{Iteracja} Kroki 2 i~3 są powtarzane aż do osiągnięcia maksymalnej wartości funkcji kosztu.
            
    \end{enumerate}

\end{tcolorbox}
